
Large language models produce poorly calibrated confidence estimates, frequently expressing high certainty on questions they answer incorrectly. This study investigates whether chain-of-thought (CoT) prompting improves calibration by enabling models to leverage in-context uncertainty signals. Through systematic mechanism verification on TruthfulQA (817 items), four steps of a hypothesized causal mechanism are empirically validated: CoT prompting produces multi-step reasoning chains (100\% rate, mean 4.49 steps), these chains contain epistemic hedging markers (82\% presence rate, mean 2.84 markers per output), markers structurally precede confidence verbalization (100\% compliance), and hedging count negatively correlates with verbalized confidence (Spearman r = -0.315, p < 1e-16, n = 664). The primary calibration claim comparing Expected Calibration Error (ECE) across conditions was not tested with real API calls due to resource constraints. These findings support a self-reading interpretation in which models integrate uncertainty signals from their own reasoning into confidence judgments, though the evidence remains correlational.
